\begin{exercice}\label{exoInf4} (\minsyndical)

Une entreprise commercialise deux produits A et B d'une m\^eme gamme : le produit A est plus cher et de meilleur qualit\'e que le B.

\medskip

Pour chaque produit, on peut mesurer sa qualit\'e, qui est la r\'ealisation d'une variable al\'eatoire $X$ : 

\begin{itemize}

\item $X$ suit une loi normale d'esp\'erance $\mu_1 = 5$ et de variance 1 si le produit est A.

\item $X$ suit une loi normale d'esp\'erance $\mu_2 = 4$ et de variance 1 si le produit est B.

\end{itemize}

\medskip

Un client ach\`ete 10 produits et veut s'assurer qu'ils sont tous du type A. Comme il vaut mieux ne pas accuser \`a tort le constructeur. Il prend comme hypoth\`ese nulle $H_0$ : $\mu = 5$ (les produits sont du type A), qu'on rejettera en prenant un risque faible. L'hypoth\`ese alternative est alors $H_1$ : $\mu = 4$. 

\medskip

On choisit comme statistique de test la moyenne du lot: $\bar{X} = \frac1n \sum_{i=1}^n X_i$.

\begin{enumerate}

\item Sous l'hypoth\`ese $H_0$, quelle est la loi de $\bar{X}$ ?

\item {\it Test 1.} On peut commencer \`a chercher un intervalle d'acceptation de la forme $A=[5-\delta ; 5 + \delta]$ tel que le risque $\alpha$ soit de $0.05$. Trouver alors la valeur de $\delta$.

\item Calculer pour ce test le risque de deuxi\`eme esp\`ece $\beta$ et la puissance du test $1-\beta$.

\item {\it Test 2.} On va maintenant construire un deuxi\`eme test, o\`u l'intervalle d'acceptation est de la forme $A = [5-\delta' ; +\infty [$. Trouver $\delta'$ pour que le risque $\alpha$ soit encore de $0.05$.

\item Calculer alors pour le Test 2 le nouveau risque de deuxi\`eme esp\`ece $\beta'$ et la puissance du test $1-\beta'$. Comparer au premier test.


\end{enumerate}

%\corrref{TD1_1}

\end{exercice}
